% JOSS-style Paper: Counter - Composable Finite-State Iteration for Python
% Based on the Journal of Open Source Software formatting

\documentclass[10pt,a4paper,onecolumn]{article}

% --- Required packages (from JOSS template) ---
\usepackage{marginnote}
\usepackage{graphicx}
\usepackage{xcolor}
\usepackage{authblk,etoolbox}
\usepackage{titlesec}
\usepackage{calc}
\usepackage{tikz}
\usepackage{hyperref}
\hypersetup{colorlinks,breaklinks=true,
            urlcolor=[rgb]{0.0, 0.5, 1.0},
            linkcolor=[rgb]{0.0, 0.5, 1.0}}
\usepackage{caption}
\usepackage{amssymb,amsmath}
\usepackage{listings}
\usepackage[T1]{fontenc}
\usepackage[utf8]{inputenc}

% --- Code listing style ---
\definecolor{codebg}{rgb}{0.95,0.95,0.95}
\lstset{
    language=Python,
    basicstyle=\ttfamily\small,
    keywordstyle=\color{blue},
    stringstyle=\color{red!70!black},
    commentstyle=\color{gray},
    numbers=none,
    frame=single,
    backgroundcolor=\color{codebg},
    breaklines=true,
    showstringspaces=false,
    columns=flexible
}

% --- Page layout (JOSS style with wide left margin) ---
\usepackage[top=3.5cm, bottom=3cm, right=1.5cm, left=1.0cm,
            headheight=2.2cm, reversemp, includemp, marginparwidth=4.5cm]{geometry}

% --- Default font (sans-serif as per JOSS) ---
\renewcommand\familydefault{\sfdefault}

% --- Section styling (JOSS style) ---
\titleformat{\section}
  {\normalfont\sffamily\Large\bfseries}
  {}{0pt}{}
\titleformat{\subsection}
  {\normalfont\sffamily\large\bfseries}
  {}{0pt}{}
\titleformat{\subsubsection}
  {\normalfont\sffamily\bfseries}
  {}{0pt}{}

% --- Header / Footer ---
\usepackage{fancyhdr}
\pagestyle{fancy}
\fancyhf{}
\renewcommand{\headrulewidth}{0pt}
\renewcommand{\footrulewidth}{0.25pt}
\fancyfoot[L]{\footnotesize\sffamily Counter: Composable Finite-State Iteration for Python}
\fancyfoot[R]{\sffamily \thepage}
\fancyheadoffset[L]{4.5cm}
\fancyfootoffset[L]{4.5cm}

% --- Author block styling ---
\renewcommand\Authfont{\sffamily\bfseries}
\renewcommand\Affilfont{\sffamily\small\mdseries}
\setlength{\affilsep}{1em}

% --- Color definitions ---
\definecolor{linky}{rgb}{0.0, 0.5, 1.0}

% --- Paragraph spacing ---
\setlength{\parindent}{0pt}
\setlength{\parskip}{6pt plus 2pt minus 1pt}

% --- Remove section numbering ---
\setcounter{secnumdepth}{0}

% --- Title and Authors ---
\title{Counter: Composable Finite-State Iteration for Python}

\author[1]{[Author Name]}
\affil[1]{[Institution], [City], [Country]}

\date{}

\begin{document}

\maketitle

% --- Sidebar (marginnote for JOSS metadata) ---
\marginpar{
  \begin{flushleft}
  \sffamily\small
  
  {\bfseries Software}
  \begin{itemize}
    \setlength\itemsep{0em}
    \item \href{https://github.com/[username]/counter}{\color{linky}{Repository}}
  \end{itemize}
  
  \vspace{2mm}
  
  {\bfseries Submitted:} December 2024
  
  \vspace{2mm}
  
  {\bfseries License}\\
  Authors of papers retain copyright and release the work under a Creative Commons Attribution 4.0 International License (\href{http://creativecommons.org/licenses/by/4.0/}{\color{linky}{CC BY 4.0}}).
  
  \end{flushleft}
}

%-------------------------------------------------------------------------------
\section{Summary}
%-------------------------------------------------------------------------------

Counter is a Python library for composable finite-state iteration with automatic 
state propagation. It provides a \texttt{Counter} class representing finite state 
spaces that can be combined using algebraic operations: product (\texttt{*}) for 
Cartesian products, sum (\texttt{+}) for disjoint unions, and synchronization for 
lockstep iteration. When iterating over a composite counter, all ancestor counters 
automatically update their state properties, eliminating the need to manually 
track indices across nested loops.

The library enforces tree-structured computation graphs, ensuring unambiguous 
state decomposition. A \texttt{CounterManager} context manager provides lifecycle 
management for counter objects, including batch reset, inactivation, and iteration 
testing utilities. Counter is designed for applications requiring systematic 
enumeration of combinatorial spaces while maintaining explicit access to component 
states.

%-------------------------------------------------------------------------------
\section{Statement of Need}
%-------------------------------------------------------------------------------

Many scientific and engineering applications require enumerating combinatorial 
state spaces: experimental designs, parameter sweeps, sequence libraries, and 
optimization landscapes. Python's standard library provides \texttt{itertools.product} 
for Cartesian products and \texttt{itertools.chain} for concatenation, while NumPy 
offers \texttt{ndindex} for multi-dimensional index iteration. However, these tools 
yield tuples or indices without maintaining persistent state objects.

Consider a common scenario: iterating over all combinations of positions and 
mutations in a biological sequence. With \texttt{itertools.product}, one must 
unpack tuples and track which component is which:

\begin{lstlisting}
for pos, mut in itertools.product(range(100), range(4)):
    process(pos, mut)
\end{lstlisting}

For deeper nesting or when different code paths need access to component states, 
this becomes unwieldy. Counter addresses this by providing stateful iteration 
where component counters expose their current state as properties:

\begin{lstlisting}
position = Counter(100, name='position')
mutation = Counter(4, name='mutation')
experiment = position * mutation

for _ in experiment:
    process(position.state, mutation.state)
\end{lstlisting}

The state propagation is automatic: advancing \texttt{experiment} updates 
\texttt{position.state} and \texttt{mutation.state} transparently.

%-------------------------------------------------------------------------------
\section{Functionality}
%-------------------------------------------------------------------------------

Counter provides the following core operations:

\begin{itemize}
    \item \textbf{Product} (\texttt{A * B}): Creates a counter over the Cartesian 
    product with $|A| \times |B|$ states. Supports N-ary products via 
    \texttt{multiply\_counters(A, B, C, ...)}.
    
    \item \textbf{Sum} (\texttt{A + B}): Creates a counter over the disjoint union 
    with $|A| + |B|$ states. During iteration, exactly one parent is active 
    (state $\geq 0$) while others are inactive (state $= -1$). Supports N-ary 
    sums via \texttt{sum\_counters(A, B, C, ...)}.
    
    \item \textbf{Synchronize} (\texttt{synchronize\_counters(A, B)}): Creates a 
    counter that keeps multiple equal-sized counters in lockstep.
    
    \item \textbf{Slice} (\texttt{A[start:stop:step]}): Creates a counter over a 
    subset of states, supporting Python slice semantics including reversal.
    
    \item \textbf{Repeat} (\texttt{A * n} or \texttt{repeat\_counter(A, n)}): 
    Creates a counter that cycles through A's states n times.
\end{itemize}

State propagation is implemented via a \texttt{decompose} method on each operation, 
which maps a composite state to parent states. For products, this uses mixed-radix 
decomposition; for sums, it identifies the active branch. The tree constraint 
(no counter appearing multiple times in a graph) ensures decomposition is 
unambiguous.

From a mathematical perspective, these operations correspond to products and 
coproducts in the category \textbf{FinSet} of finite sets, with the inactive 
state extending each set to a pointed set.

%-------------------------------------------------------------------------------
\section{Example Usage}
%-------------------------------------------------------------------------------

The following example demonstrates composing counters and observing automatic 
state propagation:

\begin{lstlisting}
from counter import Counter, CounterManager

with CounterManager() as mgr:
    # Define component counters
    row = Counter(3, name='row')
    col = Counter(4, name='col')
    
    # Compose via Cartesian product
    grid = row * col
    grid.name = 'grid'
    
    # Iterate and observe state propagation
    for state in grid:
        print(f"grid={state}, row={row.state}, col={col.state}")
\end{lstlisting}

Output:
\begin{lstlisting}[language={}]
grid=0, row=0, col=0
grid=1, row=1, col=0
grid=2, row=2, col=0
grid=3, row=0, col=1
...
grid=11, row=2, col=3
\end{lstlisting}

The \texttt{CounterManager} provides utilities for inspecting iteration behavior:

\begin{lstlisting}
df = mgr.test_iteration(grid)  # Returns DataFrame of all states
mgr.reset_all()                # Reset all counters to state 0
mgr.inactivate_all()           # Set all counters to state -1
\end{lstlisting}

%-------------------------------------------------------------------------------
\section{Acknowledgements}
%-------------------------------------------------------------------------------

[Acknowledgements to be added.]

%-------------------------------------------------------------------------------
\section{References}
%-------------------------------------------------------------------------------

\begin{enumerate}
\item Harris, C. R., et al. (2020). Array programming with NumPy. \textit{Nature}, 585(7825), 357--362. \url{https://doi.org/10.1038/s41586-020-2649-2}

\item Python Software Foundation. (2024). itertools --- Functions creating iterators for efficient looping. \url{https://docs.python.org/3/library/itertools.html}

\item Mac Lane, S. (1998). \textit{Categories for the Working Mathematician} (2nd ed.). Springer.
\end{enumerate}

\end{document}
